\section{Code Block Syntax and Indentation Semantics}

Before Python can execute branches, loops, exception handlers, functions, or classes, it must first know which statements belong together. This is the purpose of block syntax. A block is not merely a group of lines that look related to the human reader. It is a structural unit that the tokenizer and parser must recognize before the program can become executable.

Python makes an unusually direct design choice here: indentation defines block membership. The left margin is therefore not decoration. It is part of the language's structural machinery. A line moved one indentation level to the left or right may no longer belong to the same control-flow region. In Python, visual alignment is executable information.

For programmers coming from C, C++, or Java, this requires a deliberate mental adjustment. Curly-brace languages separate formal block delimiters from human formatting. Python removes that separation. The visible shape of the program and the grammatical structure of the program are forced to match.

This section explains that mechanism from both the programmer's view and the tokenizer's view: why Python uses syntactic whitespace, how indentation replaces explicit block tokens, how invalid indentation becomes a syntax error, and why mixed tabs and spaces can make the block structure ambiguous.